\documentclass{ximera}

\title{Inverse Functions Activity}
\author{MATH 425: Calculus I}

\begin{document}
\begin{abstract}
    Working with the peers in your group, solve the following problems. Make sure to show and justify all your work. Make sure everyone in the group understands the solution and participates. Be prepared to report your answers to the whole class. 
\end{abstract}
\maketitle

\begin{exercise}
     Do the functions $r(t)$ and $s(t)$ defined by
        \[r(t)=3-\frac15(t-1)^3\]
    and 
        \[s(t)=3-\frac15(t-1)^2\]
    have corresponding inverse functions? If not, use algebraic reasoning to explain why; if so, demonstrate by using algebra to find a formula for the inverse function.
    Problem from Boelkins et al., \textit{Active Prelude to Calculus}.
\end{exercise}

\begin{exercise}
    The conversion from degrees Celsius to degrees Fahrenheit is given by $F=g(C)=\frac95 C+32$.
    \begin{enumerate}
        \item Solve the equation $F=\frac95C+32$ for $C$.  We will call this function $h(C)$.
        \item Find the simplest expression that you can for the composite function $j(C)=h(g(C))$.
        \item Find the simplest expression that you can for the composite function $k(F)=g(h(F))$.
        \item Why are the functions $j$ and $k$ so simple?  Explain by discussing how $g$ and $h$ process inputs to generate outputs and what happens when we first execute one followed by the other.
    \end{enumerate}
    Problem from Boelkins et al., \textit{Active Prelude to Calculus}.
\end{exercise}


\begin{exercise}
    Find the inverse of the following functions.  Give the domain and range of the inverse. 
    \begin{enumerate}
    \item $f(x)=\frac{4x-1}{2x+3}$
    \item $g(x)=\sqrt{x^2-1}$
    \item $h(x)=\frac{2}{x^2+9}$
    \item $g(x)=\sin(3x-1)$
    \end{enumerate}
\end{exercise}




\end{document}